\documentclass[11pt]{article}

\usepackage[utf8]{inputenc}
\usepackage[margin=1in]{geometry}
\usepackage{amsmath,amssymb,mathtools,booktabs,enumitem,fancyhdr}
\usepackage[misc]{ifsym}
\usepackage{hyperref}

\setlength\textwidth{7in}
\setlength\parindent{0pt}
\oddsidemargin=-0.5cm
\textheight=665pt
\pagestyle{fancy}
\setlength{\headheight}{13.6pt}
\renewcommand{\headrulewidth}{0.1pt}
\renewcommand{\footrulewidth}{0.1pt}
\lhead{\scshape Pontificia Universidad Católica de Chile}
\rhead{\scshape IMC UC}
\cfoot{\thepage}

\newcommand{\solution}{\large\textbf{Solución}\normalsize}
\def\mat   #1{\mbox{\boldmath $\mathsf #1$}}
\def\vec   #1{\mbox{\boldmath $#1$}{}}
\def\R{\mathbb R}
\def\dx{\mbox{\(\,\mathrm{d}x\)}}
\def\ds{\mbox{\(\,\mathrm{d}s\)}}

\AtBeginDocument{\vspace*{-3.2\baselineskip}}

\begin{document}
    \begin{flushleft}
        \small
        \vspace{1pt}
        \textbf{IMT3130} \hfill
        \textbf{24/04/2026}
    \end{flushleft}
    \begin{center}
        \LARGE\textbf{Pauta Ayudantía 6}\\
        \vspace{3pt}
        \normalsize\textbf{Elementos finitos $P^1$, lifting y ensamblaje}\\
        \normalsize Ayudante: Bastián Herrera \Letter{\hspace{1pt}: \texttt{biherrera@uc.cl}}
        \rule{\textwidth}{0.05mm}
    \end{center}

\section*{Problemas y soluciones}

\begin{enumerate}[label=\textbf{\arabic*.}, leftmargin=*]

    \item \textbf{Enunciado.} En $\Omega=(0,1)$, considere una partición uniforme y el problema
    \[
    \begin{aligned}
        -\kappa u''+\sigma u&=f &&\text{en }(0,1),\\
        u(0)&=g_0,\qquad u(1)=g_1,
    \end{aligned}
    \]
    con elementos finitos $P^1$. Tome un lifting discreto, derive el sistema para los grados interiores, calcule la matriz elemental y ensamble la matriz tridiagonal reducida con las correcciones por $g_0,g_1$.

    \textbf{Solución.}
    \begin{enumerate}[label=(\alph*)]
        \item Sea $\{\varphi_i\}_{i=0}^N$ la base nodal de $P^1_h$. Un lifting discreto simple es
        \[
            G_h=g_0\varphi_0+g_1\varphi_N.
        \]
        Entonces $G_h(0)=g_0$, $G_h(1)=g_1$ y $G_h(x_i)=0$ para $i=1,\dots,N-1$. Escribimos
        \[
            u_h=u_{0,h}+G_h,
            \qquad
            u_{0,h}\in V_h^0.
        \]
        La formulación discreta es
        \[
            a(u_h,v_h)=\ell(v_h)
            \qquad\forall v_h\in V_h^0,
        \]
        donde
        \[
            a(w,v)=\int_0^1\kappa w'v'\dx+\int_0^1\sigma wv\dx,
            \qquad
            \ell(v)=\int_0^1fv\dx.
        \]
        Reemplazando $u_h=u_{0,h}+G_h$ y usando la bilinealidad de $a$,
        \[
        \begin{aligned}
            a(u_h,v_h)=\ell(v_h)
            &\Longleftrightarrow
            a(u_{0,h}+G_h,v_h)=\ell(v_h)\\
            &\Longleftrightarrow
            a(u_{0,h},v_h)+a(G_h,v_h)=\ell(v_h)\\
            &\Longleftrightarrow
            a(u_{0,h},v_h)=\ell(v_h)-a(G_h,v_h),
        \end{aligned}
        \]
        para todo $v_h\in V_h^0$.

        Ahora expandimos por separado la función con grados de libertad libres y el lifting:
        \[
            u_{0,h}=\sum_{j=1}^{N-1}U_j\varphi_j,
            \qquad
            G_h=g_0\varphi_0+g_1\varphi_N.
        \]
        Testeando con $v_h=\varphi_i$, para $i=1,\dots,N-1$, obtenemos
        \[
        \begin{aligned}
            a(u_{0,h},\varphi_i)
            &=
            \ell(\varphi_i)-a(G_h,\varphi_i)\\
            a\left(\sum_{j=1}^{N-1}U_j\varphi_j,\varphi_i\right)
            &=
            \ell(\varphi_i)-a(g_0\varphi_0+g_1\varphi_N,\varphi_i)\\
            \sum_{j=1}^{N-1}U_j\,a(\varphi_j,\varphi_i)
            &=
            \ell(\varphi_i)
            -g_0a(\varphi_0,\varphi_i)
            -g_1a(\varphi_N,\varphi_i).
        \end{aligned}
        \]
        Definiendo
        \[
            K_{ij}:=a(\varphi_j,\varphi_i),
            \qquad
            F_i:=\ell(\varphi_i),
        \]
        queda, para cada fila interior,
        \[
            \boxed{
            \sum_{j=1}^{N-1}K_{ij}U_j
            =
            F_i-g_0K_{i0}-g_1K_{iN},
            \qquad i=1,\dots,N-1.
            }
        \]
        La suma sólo recorre $j=1,\dots,N-1$ porque esos son los grados de libertad que se resuelven. Los valores de borde $U_0=g_0$ y $U_N=g_1$ ya están fijados, por lo que entran al lado derecho.

        En forma matricial, con
        \[
            I=\{1,\dots,N-1\},
            \qquad
            B=\{0,N\},
            \qquad
            \vec U_I=(U_1,\dots,U_{N-1})^\top,
            \qquad
            \vec G_B=(g_0,g_1)^\top,
        \]
        el sistema es
        \[
            \mat K_{II}\vec U_I
            =
            \vec F_I-\mat K_{IB}\vec G_B,
        \]
        Aquí $\mat K_{II}=(K_{ij})_{i,j\in I}$ contiene las interacciones entre grados libres. Por su parte, $\mat K_{IB}=(K_{ij})_{i\in I,\,j\in B}$ contiene las columnas asociadas a los nodos de borde.

        \item En $K_e=[x_e,x_{e+1}]$, las funciones base locales son
        \[
            \varphi_1^e(x)=\frac{x_{e+1}-x}{h},
            \qquad
            \varphi_2^e(x)=\frac{x-x_e}{h}.
        \]
        Sus derivadas son $-(1/h)$ y $1/h$. Para la parte difusiva,
        \[
        \begin{aligned}
            \int_{K_e}\kappa(\varphi_1^e)'(\varphi_1^e)'\dx
            &=
            \kappa\int_{x_e}^{x_{e+1}}\frac{1}{h^2}\dx
            =
            \frac{\kappa}{h},\\
            \int_{K_e}\kappa(\varphi_1^e)'(\varphi_2^e)'\dx
            &=
            \kappa\int_{x_e}^{x_{e+1}}\left(-\frac{1}{h^2}\right)\dx
            =
            -\frac{\kappa}{h}.
        \end{aligned}
        \]
        Las otras entradas se obtienen igual o por simetría. Por tanto,
        \[
            \int_{K_e}\kappa(\varphi_b^e)'(\varphi_a^e)'\dx
            =
            \frac{\kappa}{h}
            \begin{pmatrix}
                1&-1\\
                -1&1
            \end{pmatrix}_{ab},
        \]
        Para la parte de reacción,
        \[
        \begin{aligned}
            \int_{K_e}(\varphi_1^e)^2\dx
            &=
            \int_{x_e}^{x_{e+1}}
            \left(\frac{x_{e+1}-x}{h}\right)^2\dx
            =
            \frac{h}{3},\\
            \int_{K_e}\varphi_1^e\varphi_2^e\dx
            &=
            \int_{x_e}^{x_{e+1}}
            \frac{(x_{e+1}-x)(x-x_e)}{h^2}\dx
            =
            \frac{h}{6}.
        \end{aligned}
        \]
        Las entradas restantes se obtienen por simetría, y así
        \[
            \int_{K_e}\sigma\varphi_b^e\varphi_a^e\dx
            =
            \frac{\sigma h}{6}
            \begin{pmatrix}
                2&1\\
                1&2
            \end{pmatrix}_{ab}.
        \]
        Luego
        \[
            \mat K^e=
            \frac{\kappa}{h}
            \begin{pmatrix}
                1&-1\\
                -1&1
            \end{pmatrix}
            +
            \frac{\sigma h}{6}
            \begin{pmatrix}
                2&1\\
                1&2
            \end{pmatrix}.
        \]

        \item Para los grados de libertad interiores, la matriz reducida es tridiagonal:
        \[
            K_{ii}=\frac{2\kappa}{h}+\frac{2\sigma h}{3},
            \qquad
            K_{i,i+1}=K_{i+1,i}=-\frac{\kappa}{h}+\frac{\sigma h}{6}.
        \]
        En efecto, para un nodo interior $x_i$ contribuyen los elementos $[x_{i-1},x_i]$ y $[x_i,x_{i+1}]$. Cada uno aporta
        \[
            \frac{\kappa}{h}+\frac{\sigma h}{3}
        \]
        a la entrada diagonal asociada a $x_i$, luego
        \[
            K_{ii}
            =
            2\left(\frac{\kappa}{h}+\frac{\sigma h}{3}\right)
            =
            \frac{2\kappa}{h}+\frac{2\sigma h}{3}.
        \]
        La interacción entre dos nodos vecinos aparece en un único elemento, y corresponde a la entrada fuera de la diagonal de $\mat K^e$:
        \[
            K_{i,i+1}=K_{i+1,i}
            =
            -\frac{\kappa}{h}+\frac{\sigma h}{6}.
        \]
        En particular,
        \[
            \mat K_{II}=
            \begin{pmatrix}
                d&o& & \\
                o&d&o& \\
                 &\ddots&\ddots&\ddots\\
                 & &o&d
            \end{pmatrix},
            \qquad
            d=\frac{2\kappa}{h}+\frac{2\sigma h}{3},
            \qquad
            o=-\frac{\kappa}{h}+\frac{\sigma h}{6}.
        \]
        Sólo la primera y la última fila interior reciben correcciones de borde. La razón es el soporte local de las funciones nodales: $\varphi_0$ sólo se intersecta con $\varphi_1$, mientras que $\varphi_N$ sólo se intersecta con $\varphi_{N-1}$. Por tanto,
        \[
            a(\varphi_0,\varphi_i)=0\quad\text{si }i\neq1,
            \qquad
            a(\varphi_N,\varphi_i)=0\quad\text{si }i\neq N-1.
        \]
        Así, para la primera fila interior aparece el término conocido $g_0K_{10}=g_0o$:
        \[
        \begin{aligned}
            dU_1+oU_2
            &=
            F_1-og_0\\
            &=
            F_1+\left(\frac{\kappa}{h}-\frac{\sigma h}{6}\right)g_0,
        \end{aligned}
        \]
        y para la última fila interior aparece el término conocido $g_1K_{N-1,N}=g_1o$:
        \[
        \begin{aligned}
            oU_{N-2}+dU_{N-1}
            &=
            F_{N-1}-og_1\\
            &=
            F_{N-1}+\left(\frac{\kappa}{h}-\frac{\sigma h}{6}\right)g_1.
        \end{aligned}
        \]
    \end{enumerate}

    \item \textbf{Enunciado.} En el cuadrado $\Omega=(0,1)^2$ con nodos $x_1=(0,0)$, $x_2=(1,0)$, $x_3=(1,1)$ y $x_4=(0,1)$, considere los triángulos $K_1=(x_1,x_2,x_3)$ y $K_2=(x_1,x_3,x_4)$. Para
    \[
    \begin{aligned}
        -\Delta u&=1 &&\text{en }\Omega,\\
        \gamma_0u&=0 &&\text{en }\Gamma_D=\{0\}\times(0,1),\\
        \partial_nu+\beta\gamma_0u&=r &&\text{en }\Gamma_R=\{1\}\times(0,1),
    \end{aligned}
    \]
    use elementos $P^1$ conformes, construya IEN/ID/LM, ensamble las contribuciones de volumen y agregue la contribución Robin final.

    \textbf{Solución.}
    \begin{enumerate}[label=(\alph*)]
        \item La formulación discreta es: hallar $u_h\in V_h$ con $u_h(x_1)=u_h(x_4)=0$ tal que
        \[
            \int_\Omega \nabla u_h\cdot\nabla v_h\dx
            +\int_{\Gamma_R}\beta\,\gamma_0u_h\,\gamma_0v_h\ds
            =
            \int_\Omega v_h\dx+\int_{\Gamma_R}r\,\gamma_0v_h\ds
            \qquad\forall v_h\in V_h^0.
        \]
        Con numeración global $(x_1,x_2,x_3,x_4)$,
        \[
            \mathrm{IEN}(K_1)=(1,2,3),
            \qquad
            \mathrm{IEN}(K_2)=(1,3,4).
        \]
        Los nodos libres son $x_2$ y $x_3$, luego podemos usar
        \[
            \mathrm{ID}(1)=0,\qquad
            \mathrm{ID}(2)=1,\qquad
            \mathrm{ID}(3)=2,\qquad
            \mathrm{ID}(4)=0.
        \]
        Así,
        \[
            \mathrm{LM}(K_1)=(0,1,2),
            \qquad
            \mathrm{LM}(K_2)=(0,2,0).
        \]

        \item Para $K_1=(x_1,x_2,x_3)$, las funciones locales tienen gradientes
        \[
            \nabla\varphi_1=(-1,0),
            \qquad
            \nabla\varphi_2=(1,-1),
            \qquad
            \nabla\varphi_3=(0,1),
        \]
        y $|K_1|=1/2$. Como los gradientes son constantes en un elemento $P^1$,
        \[
            K^{(1)}_{ab}
            =
            \int_{K_1}\nabla\varphi_b\cdot\nabla\varphi_a\dx
            =
            |K_1|\,\nabla\varphi_b\cdot\nabla\varphi_a.
        \]
        Por ejemplo,
        \[
            K^{(1)}_{22}
            =
            \frac12(1,-1)\cdot(1,-1)
            =
            1,
            \qquad
            K^{(1)}_{23}
            =
            \frac12(0,1)\cdot(1,-1)
            =
            -\frac12.
        \]
        Calculando todas las entradas,
        \[
            \mat K^{(1)}
            =
            \frac12
            \begin{pmatrix}
                1&-1&0\\
                -1&2&-1\\
                0&-1&1
            \end{pmatrix}.
        \]
        Para $K_2=(x_1,x_3,x_4)$,
        \[
            \nabla\varphi_1=(0,-1),
            \qquad
            \nabla\varphi_2=(1,0),
            \qquad
            \nabla\varphi_3=(-1,1),
        \]
        y, usando nuevamente $K^{(2)}_{ab}=|K_2|\nabla\varphi_b\cdot\nabla\varphi_a$ con $|K_2|=1/2$,
        \[
            \mat K^{(2)}
            =
            \frac12
            \begin{pmatrix}
                1&0&-1\\
                0&1&-1\\
                -1&-1&2
            \end{pmatrix}.
        \]
        Como $f\equiv1$, en cada triángulo
        \[
            \vec F^{(e)}=\frac{|K_e|}{3}
            \begin{pmatrix}
                1\\1\\1
            \end{pmatrix}
            =
            \frac16
            \begin{pmatrix}
                1\\1\\1
            \end{pmatrix}.
        \]
        Ensamblamos sólo los grados de libertad libres $x_2,x_3$. Para $K_1$, como
        \[
            \mathrm{LM}(K_1)=(0,1,2),
        \]
        se conservan las filas y columnas locales $2$ y $3$:
        \[
            \begin{pmatrix}
                K^{(1)}_{22}&K^{(1)}_{23}\\
                K^{(1)}_{32}&K^{(1)}_{33}
            \end{pmatrix}
            =
            \begin{pmatrix}
                1&-\frac12\\[0.3ex]
                -\frac12&\frac12
            \end{pmatrix}.
        \]
        Para $K_2$, como
        \[
            \mathrm{LM}(K_2)=(0,2,0),
        \]
        el único grado libre del elemento es el nodo local $2$, que corresponde a $x_3$. Por tanto, $K_2$ sólo agrega
        \[
            K^{(2)}_{22}=\frac12
        \]
        a la entrada reducida asociada a $(x_3,x_3)$.

        Denotemos por $\mat K_{\mathrm{vol}}$ y $\vec F_{\mathrm{vol}}$ la matriz y el vector reducidos que vienen de las integrales de volumen, es decir, de los términos integrados sobre los triángulos $K_1$ y $K_2$. Esta notación todavía no incluye la contribución Robin sobre la arista $E_R$. Entonces
        \[
            \mat K_{\mathrm{vol}}=
            \begin{pmatrix}
                1&-\frac12\\[0.3ex]
                -\frac12&1
            \end{pmatrix},
            \qquad
            \vec F_{\mathrm{vol}}=
            \begin{pmatrix}
                \frac16\\[0.3ex]
                \frac13
            \end{pmatrix}.
        \]
        Para el vector de carga ocurre lo mismo. El elemento $K_1$ aporta $\frac16$ a $x_2$ y $\frac16$ a $x_3$, mientras que $K_2$ aporta $\frac16$ sólo a $x_3$. Por eso
        \[
            F_{\mathrm{vol},1}=\frac16,
            \qquad
            F_{\mathrm{vol},2}=\frac16+\frac16=\frac13.
        \]

        \item La arista Robin es $E_R=[x_2,x_3]$, de longitud $1$. Sobre esta arista, las funciones activas son las asociadas a $x_2$ y $x_3$. Parametrizamos
        \[
            (x,y)=(1,s),
            \qquad
            0\le s\le1.
        \]
        Como $\ds=ds$, las trazas de las funciones nodales libres son
        \[
            \gamma_0\varphi_{x_2}=1-s,
            \qquad
            \gamma_0\varphi_{x_3}=s.
        \]
        Entonces las entradas Robin vienen de
        \[
        \begin{aligned}
            \int_{E_R}\beta(\gamma_0\varphi_{x_2})^2\ds
            &=\beta\int_0^1(1-s)^2ds=\frac{\beta}{3},\\
            \int_{E_R}\beta(\gamma_0\varphi_{x_2})(\gamma_0\varphi_{x_3})\ds
            &=\beta\int_0^1s(1-s)ds=\frac{\beta}{6},\\
            \int_{E_R}\beta(\gamma_0\varphi_{x_3})^2\ds
            &=\beta\int_0^1s^2ds=\frac{\beta}{3}.
        \end{aligned}
        \]
        Por tanto, la contribución de matriz es
        \[
            \mat K_R=
            \frac{\beta}{6}
            \begin{pmatrix}
                2&1\\
                1&2
            \end{pmatrix},
        \]
        La contribución de carga se calcula igual:
        \[
            \int_{E_R}r\,\gamma_0\varphi_{x_2}\ds
            =
            r\int_0^1(1-s)ds
            =
            \frac r2,
            \qquad
            \int_{E_R}r\,\gamma_0\varphi_{x_3}\ds
            =
            r\int_0^1s\,ds
            =
            \frac r2.
        \]
        Así,
        \[
            \vec F_R=
            \frac{r}{2}
            \begin{pmatrix}
                1\\1
            \end{pmatrix}.
        \]
        Sumando las contribuciones de volumen y de Robin, el sistema final reducido es
        \[
            (\mat K_{\mathrm{vol}}+\mat K_R)
            \begin{pmatrix}
                U_2\\U_3
            \end{pmatrix}
            =
            \vec F_{\mathrm{vol}}+\vec F_R.
        \]
        Es decir,
        \[
            \left[
            \begin{pmatrix}
                1&-\frac12\\[0.3ex]
                -\frac12&1
            \end{pmatrix}
            +
            \frac{\beta}{6}
            \begin{pmatrix}
                2&1\\
                1&2
            \end{pmatrix}
            \right]
            \begin{pmatrix}
                U_2\\U_3
            \end{pmatrix}
            =
            \begin{pmatrix}
                \frac16\\[0.3ex]
                \frac13
            \end{pmatrix}
            +
            \frac{r}{2}
            \begin{pmatrix}
                1\\1
            \end{pmatrix}.
        \]
        Equivalentemente,
        \[
            \begin{pmatrix}
                1+\frac{\beta}{3}&-\frac12+\frac{\beta}{6}\\[0.3ex]
                -\frac12+\frac{\beta}{6}&1+\frac{\beta}{3}
            \end{pmatrix}
            \begin{pmatrix}
                U_2\\U_3
            \end{pmatrix}
            =
            \begin{pmatrix}
                \frac16+\frac r2\\[0.3ex]
                \frac13+\frac r2
            \end{pmatrix}.
        \]
        Si los datos Dirichlet en $x_1$ y $x_4$ fueran no homogéneos, esos valores no se resolverían como incógnitas. Se moverían al lado derecho mediante la corrección usual
        \[
            \vec F_I\longmapsto \vec F_I-\mat K_{IB}\vec G_B,
        \]
        incluyendo las contribuciones locales de volumen conectadas con los nodos $x_1$ y $x_4$.
    \end{enumerate}

\end{enumerate}

\end{document}
